\documentclass[10pt,a4paper]{amsart}
\usepackage{amsmath,amssymb,amsthm,graphicx}
% Retained preamble for the Stacks Project chapter snapshot in this directory.
% The chapter's own preamble.tex is a copy of the shared Stacks Project preamble; the block
% below is sliced MECHANICALLY, line for line, from it. Nothing was added, removed or
% redefined; the wrapper class, the xr-hyper/hyperref wiring and the externaldocument list
% are omitted because the assembler supplies the document class.
%
% lines taken from preamble.tex: 3, 5, 9, 11, 12, 156, 157, 158, ... (51 lines total)
%
% package coverage restored from the chapter's own preamble (2026-10-13): xy, xy-2cell

\usepackage[all]{xy}
\xyoption{2cell}
\UseAllTwocells

\newenvironment{reference}{\comment}{\endcomment}
\newenvironment{slogan}{\comment}{\endcomment}
\newenvironment{history}{\comment}{\endcomment}
\theoremstyle{plain}
\newtheorem{theorem}[subsection]{Theorem}
\newtheorem{proposition}[subsection]{Proposition}
\newtheorem{lemma}[subsection]{Lemma}
\theoremstyle{definition}
\newtheorem{definition}[subsection]{Definition}
\newtheorem{example}[subsection]{Example}
\newtheorem{exercise}[subsection]{Exercise}
\newtheorem{situation}[subsection]{Situation}
\theoremstyle{remark}
\newtheorem{remark}[subsection]{Remark}
\newtheorem{remarks}[subsection]{Remarks}
\numberwithin{equation}{subsection}
\def\lim{\mathop{\mathrm{lim}}\nolimits}
\def\colim{\mathop{\mathrm{colim}}\nolimits}
\def\Spec{\mathop{\mathrm{Spec}}}
\def\Hom{\mathop{\mathrm{Hom}}\nolimits}
\def\Ext{\mathop{\mathrm{Ext}}\nolimits}
\def\SheafHom{\mathop{\mathcal{H}\!\mathit{om}}\nolimits}
\def\SheafExt{\mathop{\mathcal{E}\!\mathit{xt}}\nolimits}
\def\Sch{\mathit{Sch}}
\def\Mor{\mathop{\mathrm{Mor}}\nolimits}
\def\Ob{\mathop{\mathrm{Ob}}\nolimits}
\def\Sh{\mathop{\mathit{Sh}}\nolimits}
\def\NL{\mathop{N\!L}\nolimits}
\def\CH{\mathop{\mathrm{CH}}\nolimits}
\def\proetale{{pro\text{-}\acute{e}tale}}
\def\etale{{\acute{e}tale}}
\def\QCoh{\mathit{QCoh}}
\def\Ker{\mathop{\mathrm{Ker}}}
\def\Im{\mathop{\mathrm{Im}}}
\def\Coker{\mathop{\mathrm{Coker}}}
\def\Coim{\mathop{\mathrm{Coim}}}
\DeclareMathSymbol{\boxtimes}{\mathbin}{AMSa}{"02}
\def\QCohstack{\mathcal{QC}\!\mathit{oh}}
\def\Cohstack{\mathcal{C}\!\mathit{oh}}
\def\Spacesstack{\mathcal{S}\!\mathit{paces}}
\def\Quotfunctor{\mathrm{Quot}}
\def\Hilbfunctor{\mathrm{Hilb}}
\def\Curvesstack{\mathcal{C}\!\mathit{urves}}
\def\Polarizedstack{\mathcal{P}\!\mathit{olarized}}
\def\Complexesstack{\mathcal{C}\!\mathit{omplexes}}
\def\Pic{\mathop{\mathrm{Pic}}\nolimits}
\def\Picardstack{\mathcal{P}\!\mathit{ic}}
\def\Picardfunctor{\mathrm{Pic}}
\def\Deformationcategory{\mathcal{D}\!\mathit{ef}}
\usepackage{cleveref}
\crefname{originallemma}{original lemma}{original lemmas}
\Crefname{originallemma}{Original lemma}{Original lemmas}
\makeatletter
\newlabel{algebra-lemma-lci-CM}{{00SB}{source}{}{sourcepointer-1}{}}
\newlabel{algebra-lemma-lci-CM@cref}{{[originallemma][0][]00SB}{[1][0][]source}{}{}{}}
\makeatother
\crefname{originallemma}{original lemma}{original lemmas}
\Crefname{originallemma}{Original lemma}{Original lemmas}
\makeatletter
\newlabel{algebra-lemma-smooth-over-field}{{00T5}{source}{}{sourcepointer-2}{}}
\newlabel{algebra-lemma-smooth-over-field@cref}{{[originallemma][0][]00T5}{[1][0][]source}{}{}{}}
\makeatother
\crefname{originallemma}{original lemma}{original lemmas}
\Crefname{originallemma}{Original lemma}{Original lemmas}
\makeatletter
\newlabel{lemma-CM-1-curves}{{0E1F}{source}{}{sourcepointer-3}{}}
\newlabel{lemma-CM-1-curves@cref}{{[originallemma][0][]0E1F}{[1][0][]source}{}{}{}}
\makeatother
\crefname{originallemma}{original lemma}{original lemmas}
\Crefname{originallemma}{Original lemma}{Original lemmas}
\makeatletter
\newlabel{spaces-morphisms-lemma-where-smooth}{{0DZI}{source}{}{sourcepointer-4}{}}
\newlabel{spaces-morphisms-lemma-where-smooth@cref}{{[originallemma][0][]0DZI}{[1][0][]source}{}{}{}}
\makeatother
\title{Stacks, Moduli of Curves: lemma-smooth-curves}
\author{The Stacks Project Authors}
\date{stacks-project:moduli-curves.tex:local-snapshot-20261012; GNU FDL 1.2 or later}
\begin{document}
\maketitle
\noindent\textit{Editorial note.} Statement span is lines 1737-1760 of sources/stacks-moduli-curves/pinned-original-source.tex; proof span is lines 1762-1787, the complete author proof (24 body lines).\par\medskip
\noindent\textit{Editorial note.} Navigation stubs: 4 cited-but-not-included labels are supplied as non-load-bearing pointers with their pinned tags where the snapshot tag index has them.\par\medskip
\noindent\textit{Editorial note.} extra\_packages is empty on purpose: the retained preamble already loads xy with options and a second declaration clashes in the assembler wrapper.\par\medskip
\noindent\textit{Editorial note.} Text is retained verbatim, including the author's own typos and placeholders.\par\medskip
\noindent\textit{Editorial note.} Queue index 17 of queue\_moduli-curves.json (pb=1140, stmt 1737-1760, proof 1762-1787).\par\medskip
\noindent\textit{Allowed standard prerequisites (editorial).}\begin{itemize}
\item Explicit prerequisite boundary. The cited inputs are prerequisites; the argument in the retained span is not replaced by any cited result.
\end{itemize}\medskip
\ifdefined\hypertarget\hypertarget{sourcepointer-1}{}\fi
\noindent\textit{Source reference.} Original lemma 00SB: sources/stacks-moduli-curves/pinned-original-source.tex, source lines 1787--1787. Cited by the retained statement or proof but not part of the collected spans; supplied only as a navigation pointer. The argument does not depend on its content.\par\medskip
\ifdefined\hypertarget\hypertarget{sourcepointer-2}{}\fi
\noindent\textit{Source reference.} Original lemma 00T5: sources/stacks-moduli-curves/pinned-original-source.tex, source lines 1787--1787. Cited by the retained statement or proof but not part of the collected spans; supplied only as a navigation pointer. The argument does not depend on its content.\par\medskip
\ifdefined\hypertarget\hypertarget{sourcepointer-3}{}\fi
\noindent\textit{Source reference.} Original lemma 0E1F: sources/stacks-moduli-curves/pinned-original-source.tex, source lines 1008--1019. Cited by the retained statement or proof but not part of the collected spans; supplied only as a navigation pointer. The argument does not depend on its content.\par\medskip
\ifdefined\hypertarget\hypertarget{sourcepointer-4}{}\fi
\noindent\textit{Source reference.} Original lemma 0DZI: sources/stacks-moduli-curves/pinned-original-source.tex, source lines 1787--1787. Cited by the retained statement or proof but not part of the collected spans; supplied only as a navigation pointer. The argument does not depend on its content.\par\medskip
\begin{lemma}
\label{lemma-smooth-curves}
There exist an open substacks
$$
\Curvesstack^{smooth, 1} \subset \Curvesstack^{smooth} \subset \Curvesstack
$$
such that
\begin{enumerate}
\item given a family of curves $f : X \to S$ the following are equivalent
\begin{enumerate}
\item the classifying morphism $S \to \Curvesstack$ factors
through $\Curvesstack^{smooth}$, resp.\ $\Curvesstack^{smooth, 1}$,
\item $f$ is smooth, resp.\ smooth of relative dimension $1$,
\end{enumerate}
\item given $X$ a scheme proper over a field $k$ with
$\dim(X) \leq 1$ the following are equivalent
\begin{enumerate}
\item the classifying morphism $\Spec(k) \to \Curvesstack$
factors through $\Curvesstack^{smooth}$, resp.\ $\Curvesstack^{smooth, 1}$,
\item $X$ is smooth over $k$, resp.\ $X$ is smooth over $k$ and
$X$ is equidimensional of dimension $1$.
\end{enumerate}
\end{enumerate}
\end{lemma}

\begin{proof}
To prove the statements regarding $\Curvesstack^{smooth}$
it suffices to show that given a family of curves
$f : X \to S$, there is an open subscheme $S' \subset S$
such that $S' \times_S X \to S'$ is smooth and such that the
formation of this open commutes with base change.
We know that there is a maximal open $U \subset X$ such
that $U \to S$ is smooth and that formation of $U$ commutes
with arbitrary base change, see
Morphisms of Spaces, Lemma \ref{spaces-morphisms-lemma-where-smooth}.
If $T = |X| \setminus |U|$ then $f(T)$ is closed in $S$ as $f$ is proper.
Setting $S' = S \setminus f(T)$ we obtain the desired open.

\medskip\noindent
Let $f : X \to S$ be a family of curves with $f$ smooth.
Then the fibres $X_s$ are smooth over $\kappa(s)$ and hence
Cohen-Macaulay (for example you can see this using
Algebra, Lemmas \ref{algebra-lemma-smooth-over-field} and
\ref{algebra-lemma-lci-CM}). Thus we see that we may set
$$
\Curvesstack^{smooth, 1} = \Curvesstack^{smooth} \cap
\Curvesstack^{CM, 1}
$$
and the desired equivalences follow from what we've already
shown for $\Curvesstack^{smooth}$ and Lemma \ref{lemma-CM-1-curves}.
\end{proof}

\end{document}
